% STATUS: written 2026-10-12
% Sources read in full: Martindale 1969 (pp. 576--578, Sections 1 and 2), Kharchenko 1991 (Sections 1.4 and 1.5, pp. 19--25), Lanski 1988 (p. 277).
\section{Martindale quotients and the extended centroid}\label{sec:prime-semiprime:quotient-rings}

The theorems of Sections~\ref{sec:prime-semiprime:posner}
and~\ref{sec:prime-semiprime:herstein-school} ended with an element $\lambda$
of the extended centroid \cite[p.~577]{Martindale1969}, which is not in the ring.
This section tells where that field comes from.

Martindale built it in 1969 \cite{Martindale1969}. He wanted one proof for a theorem of Posner on prime
rings with a polynomial identity and a theorem of Amitsur on primitive rings
\cite[p.~576]{Martindale1969}. For this he constructs, for any prime ring $R$, a
ring of quotients $Q$ whose center is a field \cite[p.~576]{Martindale1969}.

The construction uses the nonzero two-sided ideals of $R$
\cite[p.~576]{Martindale1969}. He takes all $R$-homomorphisms from such an ideal
into $R$, and calls two of them equivalent if they agree on a nonzero ideal. The
classes form a ring $Q$, and $R$ is a subring of $Q$ because $R$ is prime
\cite[p.~577]{Martindale1969}.

For each $q\in Q$ there is a nonzero ideal $U$ of $R$ with $qU\subseteq R$
\cite[p.~577]{Martindale1969}. Further, $Q$ is a prime ring
\cite[p.~577]{Martindale1969}.

Martindale shows that the center $C$ of $Q$ is a field
\cite[p.~577]{Martindale1969}. He calls it the \emph{extended centroid} of
$R$, and he calls $S=RC$ the \emph{central closure} of $R$
\cite[p.~577]{Martindale1969}.

We state the tool that the previous sections used.

\begin{theorem}[{\cite[Theorem~1, p.~577]{Martindale1969}}]\label{thm:prime:martindale-dependent}
Let $R$ be a prime ring with extended centroid $C$ and central closure $S$.
Let $a,b\in S$ be such that $axb=bxa$ for all $x\in R$. Then $a$ and $b$ are
$C$-dependent.
\end{theorem}

Herstein uses this result of Martindale \cite[p.~510]{Herstein1979}, and so
does Bre\v{s}ar \cite[p.~387]{Bresar1993}. Thus it produces the element $\lambda$ of
the earlier theorems. Then an identity $axb=bxa$ forces $b=\lambda a$ with
$\lambda$ in $C$.

Kharchenko defines the Martindale ring of quotients of a semiprime ring
\cite[Section~1.4, p.~19]{Kharchenko1991}.

Let $F$ be the set of ideals of $R$ with zero annihilator. The classes of
homomorphisms from ideals in $F$ into $R$ form the left Martindale ring of
quotients $R_F$ \cite[1.4.8, p.~21]{Kharchenko1991}.

\begin{theorem}[{\cite[Theorem~1.4.9, p.~22]{Kharchenko1991}}]\label{thm:prime:kharchenko-quotients}
Let $R$ be a semiprime ring. The ring $R_F$ has the following properties.
\begin{enumerate}
\item For every $a\in R_F$ there is an ideal $I_a\in F$ with $I_aa\subseteq R$.
\item If $a\in R_F$ and $Ia=0$ for some $I\in F$, then $a=0$.
\item If $I\in F$ and $\varphi\colon I\to R$ is a homomorphism of left
$R$-modules, then there is $r\in R_F$ with $\varphi(i)=ir$ for all $i\in I$.
\end{enumerate}
\end{theorem}

Kharchenko calls the center of $R_F$ the \emph{generalized centroid} of $R$
\cite[1.5.1, p.~24]{Kharchenko1991}. For a prime ring it is the extended
centroid of Martindale. Lanski describes $Q$ in the same way, as classes of
homomorphisms from ideals into $R$ \cite[p.~277]{Lanski1988}. He notes that the
center $C$ of $Q$ is a field, and that $C$ is also the centralizer of $R$ in $Q$.

The quotient ring matters for derivations because derivations pass to it.
Section~\ref{sec:basic-theory:extension} recorded that a derivation of a
semiprime ring that is defined on an ideal extends in a unique way to a ring of
quotients \cite[Proposition~1.8.1, p.~50]{Kharchenko1991}. Lanski states the
same for a prime ring: each derivation has a unique extension to $Q$
\cite[p.~277]{Lanski1988}.

So a derivation of $R$ can be studied inside $Q$, where the extended centroid
$C$ is a field. The question that Section~\ref{sec:prime-semiprime:kharchenko}
asks is which derivations are inner in $Q$, and what identities the values of
a derivation can satisfy.
